\documentclass[10pt,twoside]{book}
\usepackage[utf8]{inputenc}
\usepackage[T1]{fontenc}
\usepackage{geometry}
≥ometry{paper=b5, margin=1.2in, top=1.5in, bottom=1.5in}
\usepackage{microtype}
\usepackage{amsmath,amssymb}
\usepackage{graphicx}
\usepackage{xcolor}
\usepackage{titlesec}
\usepackage{titletoc}
\usepackage{fancyhdr}
\usepackage{enumitem}
\usepackage{booktabs}
\usepackage{array}
\usepackage{colortbl}
\usepackage{multirow}
\usepackage[hang]{footmisc}
\usepackage{hyperref}
ℎypersetup{
    colorlinks,
    linkcolor=black,
    citecolor=black,
    urlcolor=black,
}

% Custom colours (velvet and shadow)
\definecolor{velvetRed}{RGB}{140,30,40}
\definecolor{silkGrey}{RGB}{90,90,100}
\definecolor{inkBlack}{RGB}{20,20,20}
\definecolor{parchment}{RGB}{245,240,225}

% Page style
\pagestyle{fancy}
\fancyhf{}
\fancyhead[RE,LO]{\textit{The Velvet Touch −− A Thief's Guide to Silk and Shadow}}
\fancyhead[LE,RO]{þepage}
\fancyfoot[CE,CO]{\textcolor{silkGrey}{−−− \textit{The coin that never spends is the one you don't remember taking.} −−−}}

% Heist Engine box environment
≠wtcolorbox{heistengine}[1][]{%
  colback=black!5,%
  colframe=velvetRed!80!black,%
  title={\textcolor{velvetRed}{Heist Engine – GM Reference}},%
  fonttitle=\small\bfseries,
  sharp corners,
  enclosure title=upper left,
  #1
}

% Chapter style
\titleformat{\chapter}[display]
  {\normalfont\Huge\bfseries\color{velvetRed}}
  {\titlerule[1pt]\vspace{4pt}Łarge\chaptertitlename\ þechapter}
  {0pt}
  {}
\titlespacing*{\chapter}{0pt}{−30pt}{20pt}

\setlength{\parindent}{0pt}
\setlength{\parskip}{0.8em}

% Story / confession box
≠wenvironment{storybox}
  {\begin{quote}ℹtshape\small\setlength{\parskip}{0.5em}}
  {\end{quote}}

% Mechanics box
≠wtcolorbox{mechbox}[1]{colback=black!5,colframe=velvetRed!80!black,title={\textcolor{velvetRed}{#1}},fonttitle=\small\bfseries}

\begin{document}

% −−− TITLE PAGE −−−
\begin{titlepage}
\begin{center}
\vspace*{2cm}
{ℎuge\bfseries\color{velvetRed} The Velvet Touch}\\[0.8cm]
{Łargeℹtshape A Thief's Guide to Silk and Shadow}\\[1.5cm]
\rule{0.6\textwidth}{1pt}\\[0.5cm]
{łarge Madam Serafine, Mistress of the Velvet Court}\\[0.3cm]
{\em Silkstrand, Year of the Falling Bridge}\\[1cm]
\vfill
{\small \textcolor{silkGrey}{The coin you take is never the one that costs you. The one you leave behind−−−that is the debt that collects.}}
\end{center}
\end{titlepage}

% −−− FRONT MATTER −−−
\chapter*{Preface: On Rising from the Gutter}
\addcontentsline{toc}{chapter}{Preface}

I was seven when I picked my first pocket. A merchant from Kahfagia, fat with spices and arrogance. He never felt my fingers slip his purse. I watched him laugh with his friends while I bought bread for a week.

They never caught me. Not because I was clever−−−I was not, not yet. Because I was \emph{nothing}. A gutter rat, a ghost in the crowd, a face that blends into wet cobblestones. That is the first lesson: power is not in the blade or the coin. It is in being \emph{forgotten}.

Twenty years later, I own three warehouses, two gambling dens, and a brothel that also launders secrets. The Watch Captain owes me his promotion. The Dye Syndicate pays me not to undercut their prices. And the Archivolt notaries keep a drawer full of my forged seals−−−they do not know which are real.

I am Madam Serafine. They call me the Velvet Touch because I never leave bruises.

This book is not for heroes. It is for those who understand that the world runs on whispers, debts, and the careful omission of the truth. I learned from the streets. I learned from Ikasha, who sleeps in shadow, and from Livaea, who dresses power in silk and smiles. I learned that a song can open a lock faster than any pick, and that a forged document is worth a thousand swords.

You want to be a thief? Good. But first, forget everything you think you know. The mark is not the victim. The mark is the one who believes they are safe.

The coin that never spends is the one you don't remember taking.

\vspace{0.5cm}
\begin{flushright}
Madam Serafine\\
Silkstrand, under the Three−Queens Bridge\\
The ink is lamp−black and my blood. The mirror is clean.
\end{flushright}

\tableofcontents
≠wpage

% −−−−−−−−−−−−−−−−−−−−−−−−−−−−−−−−−−−−−−−−−−−−−−−−−−−−−−−−−−−−−
% PART 1: THE PHILOSOPHY OF THE QUIET HAND
% −−−−−−−−−−−−−−−−−−−−−−−−−−−−−−−−−−−−−−−−−−−−−−−−−−−−−−−−−−−−−
\part{The Philosophy of the Quiet Hand}

\chapter{On Being Nothing}

The first time you break into a counting−house, you will be afraid. Good. Fear keeps your fingers light and your ears open. The moment you feel safe, you are already caught.

A thief is not a warrior. We do not fight. We \emph{slip}. Through a window left unlatched, through a crowd that parts without knowing, through a contract signed in too much haste. The world is full of gaps−−−the ninth step un−counted, the witness who looked away, the seal that was pressed a moment too soon. 

\section{The Three Principles of the Velvet Hand}

\subsection{The Coin That Is Not There}

Never take what will be missed before dawn. A silver coin from a rich man's purse? He will curse his servants, forget by noon. A ledger page that proves his fraud? He will burn the city to find it.

Take what they do not know they have \emph{until it is too late}.

\subsection{The Witness Who Looks Away}

Every crime is witnessed. The trick is not to be invisible−−−it is to be \emph{unremarkable}. A servant carrying linens. A clerk with a stack of writs. A beggar sleeping in the archway. The guards see you. They just do not \emph{notice}.

\subsection{The Debt That Pays Twice}

When you steal, you create a debt. Not yours−−−\emph{theirs}. They owe you for the fear you planted, the reputation you damaged, the secret you now hold. A true thief does not take gold. Gold spends. A secret spends forever.

\chapter{Thief Talents: The Velvet Ladder}
łabel{ch:thief−talents}

The following talents are suited for those who walk the edge of law and shadow. They are not for brawlers or brutes. They are for the quiet, the cunning, and the patient.

\section{2 XP Talents −−− Foundations}

\begin{talentbox}{Light Fingers}
\textbf{Tags:} Strike, AGI, Gambit \\
\textbf{Cost: 2 XP} \\
\textbf{Active, once per scene.} When you attempt to lift a small object (purse, key, letter, seal) from a person or desk, treat your first Stealth or Subterfuge roll as Position +1.

No one notices until it is already in your sleeve. (See \ref{talent:light−fingers})
\end{talentbox}

\begin{talentbox}{The Unlocked Window}
\textbf{Tags:} Move, AGI, Flow \\
\textbf{Cost: 2 XP} \\
\textbf{Passive.} When infiltrating a structure you have cased for at least an hour, ignore one security obstacle (locked door, patrolled corridor, guard dog) per session. Describe how you already prepared for it. (See \ref{talent:unlocked−window})
\end{talentbox}

\begin{talentbox}{Silk Palms}
\textbf{Tags:} Influence, FAT \\
\textbf{Cost: 2 XP} \\
\textbf{Active.} When you bribe or gift a minor official, you gain a temporary String: \textit{Bought Silence}. Once, you may activate this String to reduce DV by 1 on a Petition or Broker action involving that official. The String is consumed afterward. (See \ref{talent:silk−palms})
\end{talentbox}

\section{4 XP Talents −−− The Working Knife}

\begin{talentbox}{Always Another Exit}
\textbf{Tags:} Move, SPT, Gambit \\
\textbf{Cost: 4 XP} \\
\textbf{Active, once per session.} When you are cornered or a heist goes wrong, you may declare you prepared an escape route earlier. Describe it plausibly (rope out a window, bribed a guard, secret door). The GM may impose a cost−−−a dropped item, a witness, a future debt−−−but you \emph{get out}. (See \ref{talent:always−another−exit})
\end{talentbox}

\begin{talentbox}{The Forger's Eye}
\textbf{Tags:} Craft, WIT, Timer \\
\textbf{Cost: 4 XP} \\
\textbf{Active.} You can copy any handwriting, seal, or signature with uncanny accuracy. When you attempt to forge a document, roll Wits + Craft (DV 3 for simple, 4−−5 for complex). On a hit, create a \textbf{Forged credential}. On a partial, the forgery holds for one scene; thereafter it may be detected. (See \ref{talent:forgers−eye}) \\
\textit{Requires: Craft 1+.}
\end{talentbox}

\begin{talentbox}{Crowd Veil}
\textbf{Tags:} Observe, MIG, Support \\
\textbf{Cost: 4 XP} \\
\textbf{Active.} You vanish into crowds. When in a populated area (market, festival, dock), you may roll Stealth vs. DV 3 to become \textit{Hidden in plain sight}. Pursuers suffer −−1 die to find you until you act conspicuously. (See \ref{talent:crowd−veil}) \\
\textit{Requires: Stealth 2+.}
\end{talentbox}

\section{5 XP Talents −−− Signature Moves}

\begin{talentbox}{The Velvet Glove}
\textbf{Tags:} Influence, Strike, Combo \\
\textbf{Cost: 5 XP} \\
\textbf{Active, once per session.} You have cultivated a network of informants, fence contacts, and corrupted clerks. You may declare that a minor NPC (a servant, a dockworker, a scribe) owes you a favor. They provide a specific service: unlock a door, delay a patrol, deliver a message. The GM may dictate a cost (the NPC will later need help, or their favor expires after use). (See \ref{talent:velvet−glove}) \\
\textit{Requires: Presence 2+, Streetwise 1+.}
\end{talentbox}

\begin{talentbox}{Master of Misdirection}
\textbf{Tags:} Influence, Gambit, SB \\
\textbf{Cost: 5 XP} \\
\textbf{Active.} When you are discovered during a crime, you may spend 1 Boon to redirect suspicion to another person or faction. Roll Deception vs. DV 4. On a hit, the heat shifts; on a miss, you gain Exposure +1 and the target becomes hostile. (See \ref{talent:master−of−misdirection}) \\
\textit{Requires: Deception 2+.}
\end{talentbox}

\section{7 XP Prestige Talent −−− The Quiet Catastrophe}

\begin{talentbox}{Ghost Heist}
\textbf{Tags:} Strike, Timer , Overload \\
\textbf{Cost: 7 XP} \\
\textbf{Active, once per arc.} You may commit a crime that is never solved−−−not because you are perfect, but because the evidence points nowhere. After a successful heist, you may declare that no traceable evidence remains. The GM cannot use Social SB to retroactively implicate you from that scene. \textbf{Cost:} You must accept a permanent \textit{Reputation: Whispers} tag; some will always suspect you, even without proof. (See \ref{talent:ghost−heist}) \\
\textit{Requires: Tier II, Stealth 3+.}
\end{talentbox}

\chapter{Rogue Builds for the Silkstrand Night}

\section{The Cat Burglar}
\textbf{Core Talents:} Light Fingers, The Unlocked Window, Always Another Exit.  
\textbf{Style:} High, vertical, through windows and rooftops.  
\textbf{Patron:} Ikasha.

\section{The Confidence Artist}
\textbf{Core Talents:} Silk Palms, The Forger's Eye, Velvet Glove.  
\textbf{Style:} Social infiltration, false identities, forged documents.  
\textbf{Patron:} Livaea.

\section{The Saboteur}
\textbf{Core Talents:} Crowd Veil, Master of Misdirection, The Ghost Did It.  
\textbf{Style:} Disruption, frame−ups, political chaos.  
\textbf{Patron:} Malachai (rare) or Aveh.

% −−−−−−−−−−−−−−−−−−−−−−−−−−−−−−−−−−−−−−−−−−−−−−−−−−−−−−−−−−−−−
% PART 2: THE FORGER'S ART
% −−−−−−−−−−−−−−−−−−−−−−−−−−−−−−−−−−−−−−−−−−−−−−−−−−−−−−−−−−−−−
\part{The Forger's Art}

\chapter{The Weight of Paper}

In Silkstrand, the right seal is worth more than a ship. A stamped manifest can move contraband past the harbor master. A notary's ribbon can turn stolen goods into legitimate auction lots. A letter of credit can buy a warehouse with no coin changing hands.

Forgery is not crime. It is \emph{reality editing}.

\section{The Forging Mini−Game}

To forge a document, you need three things:

\begin{enumerate}
ℹtem \textbf{Knowledge} of what the genuine document looks like−−−its paper, ink, seal, and language.
ℹtem \textbf{A Sample} −−− at least one authentic example to copy, or a witness who can describe it in detail.
ℹtem \textbf{A Workshop} −−− a quiet place with the right tools: quills, inks, sealing wax, a press, and a steady hand.
\end{enumerate}

\subsection{The Four Steps of Forgery}

\paragraph{Step 1: Research} (Wits + Investigation, DV 2−−4)  
Learn the document's specifics. On a hit, gain +1 die to the Craft roll. On a miss, you miss a crucial detail−−−the forgery will be detectable.

\paragraph{Step 2: Materials} (Wits + Streetwise, DV 3)  
Source matching paper, ink, and seal. On a hit, the materials are adequate. On a partial, they are close but flawed (the forgery holds for one scene, then degrades). On a miss, you must substitute inferior materials−−−the forgery starts with a \textit{Suspicious} tag.

\paragraph{Step 3: The Craft} (Wits + Craft, DV 3 for simple (receipt, note), 4 for standard (manifest, writ), 5 for complex (royal charter, embassy seal))  
Roll to create the forged document. On a hit, you produce a \textbf{Forged Diamond}. On a partial, produce a \textbf{Fragile Diamond}−−−it works once, then crumbles (detected). On a miss, the forgery is worthless; you may salvage materials for another attempt, but the GM gains 1 SB.

\paragraph{Step 4: Presentation} (Presence + Deception vs. DV of the examiner)  
When you present the forgery, roll against the target's Insight or relevant skill. On a hit, the document is accepted. On a partial, they are suspicious−−−offer concession (bribe, favor, or they keep the document for "verification"). On a miss, they detect the forgery; your Exposure ticks +1 and the document is seized.

\subsection{Forged Diamond Types}

\begin{itemize}
ℹtem \textbf{Cargo Manifest} −− Legal immunity for smuggling. Lasts one voyage.  
ℹtem \textbf{Notarial Seal} −− Validates a contract or deed. Lasts one arbitration.  
ℹtem \textbf{Letter of Credit} −− Withdraw funds from a named merchant house. One use.  
ℹtem \textbf{Writ of Passage} −− Skip tolls and inspections on a named road or river. One journey.  
ℹtem \textbf{Identity Papers} −− Assume a false identity for one scene.  
ℹtem \textbf{Pardon} −− Commutes a minor crime. Requires a public reading.  
\end{itemize}

\subsection{Failure and Consequences}

When a forgery is detected, the GM may choose one:

\begin{itemize}
ℹtem \textbf{Exposure} −− Your reputation among officials drops. Tick Exposure [4] for the city or faction.
ℹtem \textbf{Debt} −− The examiner does not report you−−−instead, they demand a favor.
ℹtem \textbf{Bounty} −− A warrant is quietly issued. A rival thief may claim it.
ℹtem \textbf{Counterfeit Trail} −− Your forged seal is used by someone else later, framing you.
\end{itemize}

\chapter{Fence and Laundry}

Once you have stolen goods, you need to sell them. The art of laundering is not hiding−−−it is \emph{recontextualizing}.

\subsection{Market Tiers}
Different fences operate at different levels of legitimacy and scrutiny, affecting the difficulty of moving goods:
\begin{itemize}
ℹtem \textbf{Street Market (DV 2)}: Stolen goods sold openly; quick but risky (e.g., pawn shops, street vendors)
ℹtem \textbf{Guild Market (DV 3)}: Hot cargo requiring some paperwork; moderate risk (e.g., guild merchants, legitimate fronts)
ℹtem \textbf{Court Market (DV 4)}: Tracked items needing forged documents; lower risk but higher cost (e.g., nobles, officials)
ℹtem \textbf{Arcane Market (DV 5)}: Warded or psionically significant items; highest risk and reward (rare, supernatural goods)
\end{itemize}
Partial results when fencing: price spike (Bank −1), Crew Heat +1, quality flaw (−1 Effect on next use), or seller's temporary String.
On a complete miss: stingspring (Sting[6]+2), counterfeit gear, or owed favor to a rival fence.


\subsection{Laundering Timers}

To launder a stolen item or forged Diamond, start a \textbf{Laundry Timer } [4] (small items) or [6] (major items). Each downtime scene, you may work one tick by:

\begin{itemize}
ℹtem Bribing a fence (spend 1 Favor or 1 Boon)
ℹtem Moving goods through a neutral market (risk Exposure +1)
ℹtem Using a genuine paper trail (requires a Forged Diamond)
\end{itemize}

When the timer fills, the item is clean−−−it can be sold at full value or used without suspicion.

If you fail to complete the timer within three sessions, the item becomes \textit{Hot}. Selling it then risks Exposure +2 and a visit from the original owner's agents.

% −−−−−−−−−−−−−−−−−−−−−−−−−−−−−−−−−−−−−−−−−−−−−−−−−−−−−−−−−−−−−
% PART 3: THE CANTOR'S CADENCE
% −−−−−−−−−−−−−−−−−−−−−−−−−−−−−−−−−−−−−−−−−−−−−−−−−−−−−−−−−−−−−
\part{The Cantor's Cadence}

\chapter{The Voice That Opens Doors}

Singing is not for taverns alone. A well−timed tune can calm a crowd, distract a guard, or deliver a secret message in plain hearing. The Cantor's path is the thief's second skin−−−especially when the singer cares more for the lock than the lyric.

\section{The Rogue Cantor: Song as Setup}

Most Cantors perform for applause. A rogue Cantor performs for \emph{opportunity}. While the audience watches the singer, their pockets are already empty.

\subsection{Performance as Distraction}

When you sing or play in a public space, roll Presence + Performance against a DV set by crowd size (2 for small, 4 for large). On a hit, you create an \textbf{Engrossed Audience} tag. For the duration of your song, any ally attempting Stealth or Subterfuge in the same area gains +1 die.

On a partial, the crowd is attentive but some watch the exits−−−Stealth rolls are unchanged, but allies gain +1 Position to social rolls.

\subsection{The Encore Gambit}

Once per session, after a successful performance, you may declare an \textit{Encore}. The crowd demands more. This extends your distraction for an additional scene, but the GM gains 1 SB (the guards are curious now).

\section{Corrupted Cantors: The Cult of the Velvet Note}

Not all Cantors walk the road alone. Some find patrons in the shadows−−−beings who grant power in exchange for a different kind of song.

\subsection{Ikasha's Whisper (Shadow Cantor)}

Ikasha speaks in the space between notes. Followers learn to sing in frequencies that do not carry, to hum a lullaby that makes guards sleep, to whistle a tune that unlocks spirit−locks.

\textbf{Corruption:} Your voice develops a second timber−−−a whisper that follows your sung words. Strangers shiver when you speak. Allies sometimes hear you when you are not there.

\textbf{Rites:}
\begin{itemize}
ℹtem \textit{Cradle Song} (Low) −− Lull a single target (Resist DV 3). Costs 1 Fatigue.
ℹtem \textit{Lockpick's Refrain} (Standard) −− Unlock one mundane or warded lock by humming its true pitch. Costs 1 Obligation.
\end{itemize}

\subsection{Livaea's Velvet Hook (Social Cantor)}

Livaea is the patron of courtiers, concubines, escorts, and companions−−−all those who trade in desire. She teaches that every song is a seduction, every glance a contract, every touch a negotiation. Her followers do not charm; they \emph{bind}. A promise whispered in her name is a thread that tightens with every breath.

\textbf{Corruption:} Your eyes take on a reflective quality. In dim light, others see their own desires mirrored there. Your skin always feels faintly warm to the touch. Those who sleep near you dream of lost lovers.

\textbf{Rites:}
\begin{itemize}
ℹtem \textit{Golden Tongue} (Low) −− Gain +2 dice to Sway for one social exchange. Costs 1 Fatigue; after use, the target feels a mild attraction or obligation. They may seek you out later.
ℹtem \textit{The Velvet Invitation} (Low) −− Gain automatic entry to any private social function (masquerade, salon, pleasure house). The invitation appears in your pocket, signed by a name you can borrow for one night. Costs 1 Obligation.
ℹtem \textit{The Unrefusable Offer} (Standard) −− Sing or whisper a bargain while touching the target (a hand on the arm, a brush of fingers). The target must accept the offer or suffer −2 dice to all rolls until they do. The offer must be something they could conceivably agree to (no suicide, no self−mutilation). Costs 2 Obligation. After use, the target remembers the moment as strangely intimate.
ℹtem \textit{The Crimson Masquerade} (High) −− Assume the appearance of any person you have touched (skin−to−skin) for one scene. Your voice matches theirs perfectly. Costs 2 Fatigue and marks a \textit{Debt Timer } [4] with Livaea.
\end{itemize}

\textbf{Societal Impact:} Livaea's cult is woven into every pleasure house, every high−class escort service, every discreet companionship agency in the Amaranthine. Her temples are not buildings; they are the back rooms of brothels, the upper floors of massage houses, the private salons of courtesans who buy and sell secrets along with flesh. A Livaean courtesan is as dangerous as any assassin−−−her weapons are a smile, a tear, a carefully timed yawn. The Velvet Court has a standing agreement with the Livaean network: information flows both ways, and a thief who insults a Livaean courtesan finds every door closed and every guard alerted.

\subsection{Malachai's False Note (Desperate Cantor)}

Malachai answers when the thief has nothing left to lose. Her Cantors sing of easy scores, of curses as blessings, of the coin that spends itself. The price is always a piece of the singer's soul. She is the patron of gamblers, addicts, and those who have sold themselves one too many times.

\textbf{Corruption:} Your voice cracks on certain notes. Those who hear you feel a brief, inexplicable hunger. Your shadow seems thinner, hungrier. You find yourself counting coins you never had.

\textbf{Rites:}
\begin{itemize}
ℹtem \textit{The Lucky Pick} (Low) −− Reroll a failed Stealth or Subterfuge roll. Costs 1 Fatigue and 1 Corruption tick.
ℹtem \textit{The Debt Note} (Standard) −− Sing a curse on a rival. They suffer −1 die on their next heist. Costs 2 Fatigue and marks a \textit{Debt Timer } [4] that Malachai will collect. The collection is always painful and often ironic.
ℹtem \textit{The Final Score} (High) −− Once per campaign, you may automatically succeed at any single heist roll, but you must immediately mark 2 Corruption and gain a permanent \textit{Marked by Malachai} condition. The coin you take will never be enough.
\end{itemize}

\chapter{Building a Rogue Cantor}

\section{The Whisper (Shadow Focus)}
Talents: Ikasha's rites, Stealth, Performance.  
Role: Infiltrator, messenger, spy.  
Progression: Low rites → Crowd Veil → Ikasha's Mark.

\section{The Siren (Social Focus)}
Talents: Livaea's rites, Sway, Silk Palms.  
Role: Confidence artist, blackmailer, courtesan−spy.  
Progression: Golden Tongue → The Velvet Invitation → The Unrefusable Offer → Velvet Glove.

\section{The Courtesan (Seduction Focus)}
Talents: Livaea's rites, Silk Palms, The Velvet Glove.  
Role: Escort, companion, concubine, information broker.  
Progression: The Velvet Invitation → Golden Tongue → The Unrefusable Offer → The Crimson Masquerade.

\section{The Cornered Rat (Desperate Focus)}
Talents: Malachai's rites, Always Another Exit, Master of Misdirection.  
Role: Saboteur, desperate fugitive, self−destructive genius.  
Progression: The Lucky Pick → Blood Price → The Final Score.

% −−−−−−−−−−−−−−−−−−−−−−−−−−−−−−−−−−−−−−−−−−−−−−−−−−−−−−−−−−−−−
% PART 4: THE VELVET COURT (Expanded)
% −−−−−−−−−−−−−−−−−−−−−−−−−−−−−−−−−−−−−−−−−−−−−−−−−−−−−−−−−−−−−
\part{The Velvet Court}

\chapter{The Court That Never Sleeps}

The Velvet Court is not a building. It is a network−−−of whispers, of debts, of favours owed and collected. It has no single leader, though many claim the title. It has no charter, no treasury, no walls. It is the space between the legitimate and the condemned, and its members are the ones who know how to walk that line.

I am Madam Serafine, and I hold the largest share of the Court's influence. But I did not build it alone. No one builds a shadow. You simply learn to stand where it falls.

\section{The Hierarchy of Shadows}

The Velvet Court is not a guild. There are no titles, no uniforms, no initiation rituals. Yet every thief in Silkstrand knows who to ask, who to fear, and who to pay. The Court's structure is a web of \textbf{Creditors} and \textbf{Debtors}. Reputation is currency. Favours are interest.

\subsection{The Creditors (The Velvet Seats)}

Six individuals hold the majority of the Court's influence. Each controls a different resource:

\begin{enumerate}
ℹtem \textbf{Madam Serafine} (Information, Forgery, Laundering)
ℹtem \textbf{Old Kes} (Fence, Smuggling, Markets)
ℹtem \textbf{The Dockmaster} (Movement, Ships, Hidden Warehouses)
ℹtem \textbf{Sister Agatha} (Cover, Sanctuary, False Identities)
ℹtem \textbf{The Archivist} (Legal Precedent, Forged Writs, Court Delays)
ℹtem \textbf{The Grey Eminence} (Unknown; rumoured to be a Veiled Sister or a Covenant judge)
\end{enumerate}

Each Creditor has a \textbf{Favour Timer } [6]. When you do a significant service for them, tick the timer. When it fills, they owe you a major favour (a heist's worth of help). When you owe them a debt, they may call it in at any time.

\subsection{The Debtors (Everyone Else)}

Every thief, fence, smuggler, and corrupt clerk in Silkstrand owes someone. The Court tracks debts through \textbf{Strings}. A String can be:

\begin{itemize}
ℹtem \textit{Intelligence} (a secret, a map, a guard schedule)
ℹtem \textit{Safe Passage} (a route, a blind eye, a forgotten warrant)
ℹtem \textit{Muscle} (a favour from a bruiser, a crew for a job)
ℹtem \textit{Coin} (silver, goods, or a Forged Diamond)
\end{itemize}

Strings are traded, sold, and called in. The Court's ledger is oral−−−written ink can be burned. A remembered debt is forever.

\section{Faction Reputation: The Velvet Dial}

Track the party's standing with the Velvet Court using a \textbf{Velvet Dial} from −3 (Enemy of the Court) to +3 (Creditor−in−Waiting). Adjust the dial when:

\begin{itemize}
ℹtem Complete a job for a Creditor (+1)
ℹtem Pay off a major debt (+1)
ℹtem Betray a fellow Court member (−1)
ℹtem Get caught by the Watch and name names (−2)
ℹtem Kill a Creditor (reset to −3, but gain a permanent enemy)
\end{itemize}

\textbf{Effects:}
\begin{itemize}
ℹtem \textbf{−3 to −2: Hunted.} Guards are tipped off; fences refuse you; you cannot find a safe house in Silkstrand.
ℹtem \textbf{−1 to +1: Neutral.} Standard access to minor fences, but Creditors watch you warily.
ℹtem \textbf{+2 to +3: Trusted.} You may request a Favour from any Creditor once per arc; you may attend the Velvet Moot.
\end{itemize}

\subsection{The Velvet Moot}

Once per season, the Creditors meet at a secret location. If you have Velvet Dial +2 or higher, you may be invited. At the Moot, you may:

\begin{itemize}
ℹtem \textbf{Bid on a Job} (gain exclusive access to a high−value heist)
ℹtem \textbf{Settle a Debt} (pay off a major obligation with a rare item or secret)
ℹtem \textbf{Propose a New Creditor} (challenge a current Creditor for their seat−−−risky)
ℹtem \textbf{Learn Rumours} (gain 1d3 Strings of intelligence about future events)
\end{itemize}

\chapter{The Court's Assets (What the GM Can Offer)}

When players work for the Velvet Court, they can earn access to:

\begin{itemize}
ℹtem \textbf{The Blue Ledger} (a record of every bribeable official in Silkstrand)
ℹtem \textbf{The Maze of Whispers} (a network of hidden passages under the city)
ℹtem \textbf{The Ghost Crew} (a team of urchin spies, all Cap 1, good for reconnaissance)
ℹtem \textbf{The Scribe's Seal} (a notary who will stamp anything for a price)
ℹtem \textbf{The Unstable Bridge} (a safe house that can be collapsed to cover an escape)
\end{itemize}

\subsection{Velvet Heat Dial}
The Court tracks two timers that measure active attention from authorities and rivals:
\begin{itemize}
ℹtem \textit{Watch Interest [6]} (Crew Heat) – fills when jobs go loud or evidence left.
ℹtem \textit{Syndicate Pressure [6]} (City Heat) – fills when rivals or the Exchange take notice.
\end{itemize}
At 3 segments: \textit{Patrol Sweep} begins (guards actively look).
At 6 segments: \textit{Sting Operation} triggers.

\noindent\textbf{Watch Interest [6]:} 0◉◉◉◉◉◉
\begin{itemize}
ℹtem 0−2: Background hum.
ℹtem 3: Patrol Sweep [4] begins.
ℹtem 4−5: Random stops; Notoriety +1 on flashy moves.
ℹtem 6: Sting [6] + Exposure +1
\end{itemize}

\noindent\textbf{Syndicate Pressure [6]:} 0◉◉◉◉◉◉
\begin{itemize}
ℹtem 0−2: Routine.
ℹtem 3: Curfew murmurs, stop and search; start Patrol Sweep [4] in hot districts.
ℹtem 6: Crackdown [4] (DV +1 on criminal actions).
\end{itemize}

\chapter{Jobs from the Velvet Court (d12)}

When the players need work (or the GM needs a hook), roll or choose:

\begin{enumerate}
ℹtem \textbf{The Unopened Letter} −− Steal a sealed diplomatic pouch from a courier before dawn. The contents are a treaty that would spark a trade war. Rival thieves are also hunting it.
ℹtem \textbf{The Ghost Ship} −− A vessel in the harbour carries a cargo of cursed idols. Smuggle them off without alerting the crew, who are cultists.
ℹtem \textbf{The Silent Ledger} −− A merchant's private ledger records debts to the Velvet Court. He is threatening to go to the Watch. Steal the ledger and replace it with a forgery.
ℹtem \textbf{The Fencing Master} −− A duellist is blackmailing a Creditor. Challenge him to a public duel and make him lose. No lethal damage−−−just humiliation.
ℹtem \textbf{The Last Heir} −− A noble's lost son has been located in a debtor's prison. Break him out. The noble will pay in land, not coin.
ℹtem \textbf{The Poisoned Well} −− Someone is poisoning the Court's informants. Find the source. It may be a rogue alchemist in the Dye District.
ℹtem \textbf{The Crimson Seal} −− A forged notarial seal is being sold to rival thieves. Steal the original die and destroy all copies. The forger is a former Court member.
ℹtem \textbf{The Velvet Trap} −− A new Creditor has appeared, promising to sell the Court's secrets to the Watch. Prove they are a plant, or eliminate them quietly.
ℹtem \textbf{The Night Market Heist} −− The Dust Market is auctioning an item the Court wants. Steal it before the gavel falls. The auctioneer is a rival faction's fence.
ℹtem \textbf{The Ghost of the Archivolt} −− A dead notary's ghost is haunting the court records. It is trying to expose old forgeries. Lay it to rest or silence it.
ℹtem \textbf{The Broken Bridge} −− The Three−Queens Bridge is being repaired. The construction crew has found a hidden vault beneath the old foundation. Empty it before the city engineers arrive.
ℹtem \textbf{The Ninth Key} −− A legend says the Velvet Court's founder hid a master key that can open any lock in Silkstrand. The key is split into nine fragments, each held by a different thief. The party is tasked with collecting them−−−or stealing them from the current holders.
\end{enumerate}

\chapter{Notable Court NPCs}

\subsection{Madam Serafine (Creditor)}
\textbf{Tier:} IV  
\textbf{Motivation:} Maintain the Court's balance, expand her influence to the Brass Coast.  
\textbf{Leverage:} She holds information on every major official in Silkstrand. She also knows the secret passage into the Archivolt's sealed vault.  
\textbf{Quirk:} She never wears the same dress twice. Each one has a hidden pocket large enough for a ledger.

\subsection{Old Kes (Fence)}
\textbf{Tier:} III  
\textbf{Motivation:} Increase his share of the smuggling trade, undermine the Kahfagian factors.  
\textbf{Leverage:} He can move any good, anywhere, within a week. He also has a contract with a Sidhi ship captain that is worth a small fortune.  
\textbf{Quirk:} He is missing three fingers on his left hand. He lost them to a customs trap. He keeps them in a jar as a warning.

\subsection{Sister Agatha (Sanctuary)}
\textbf{Tier:} II (but with powerful allies)  
\textbf{Motivation:} Protect the poor and the desperate. She gives sanctuary to thieves on the run, but demands they work off their stay in her soup kitchen.  
\textbf{Leverage:} Her hospice is neutral ground. No violence is permitted within its walls. The Watch respects her because she heals their wounded.  
\textbf{Quirk:} She has a tattoo of a broken chain on her wrist. She was a freed slave from Ashaan.

\subsection{The Archivist}
\textbf{Tier:} III  
\textbf{Motivation:} Accumulate legal knowledge and use it to blackmail officials. He is also the secret keeper of the Court's forged seal database.  
\textbf{Leverage:} He can make any legal problem disappear−−−or appear. He also knows which notaries are corrupt and which are clean.  
\textbf{Quirk:} He has never been seen without a specific quill, a family heirloom. The quill's nib is made from a saint's finger bone.

\subsection{The Grey Eminence (Secret)}
\textbf{Tier:} ? (rumoured IV)  
\textbf{Motivation:} Unknown. Some say they are a Veiled Sister who fled Ashaan. Others say they are a Covenant judge playing a long game.  
\textbf{Leverage:} The Grey Eminence is never seen, never named, but their orders carry weight. They may be the only one who knows the true origin of the Velvet Court.  
\textbf{Quirk:} Letters from the Grey Eminence arrive on paper that is still warm, as if freshly pressed.

\subsection{Informant Web Mechanics}
The party can cultivate 3−5 informants (dock porter, lamplighter, clerk, etc.). Each has Reliability [4] (ticks down on 1s). Pay their Price (a favor, a coin, a secret) for a Clue or DV −1 on one roll. Two informants on the same topic: +1 Effect, only lowest Reliability ticks down.

\noindent\textbf{Running a Web}
\begin{itemize}
ℹtem \textbf{Access}: What information the informant can provide
ℹtem \textbf{Reliability [4]}: Tracks trustworthiness; decreases when the informant is stressed or compromised (decreases on a roll of 1 when used)
ℹtem \textbf{Price}: Cost to gain information (favor, coin, secret)
ℹtem \textbf{Flags}: Special traits (e.g., \"Connected to Watch\", \"Lives in Dye District\")
\end{itemize}

\noindent\textbf{Using Informants}
Pay Price to gain:
\begin{itemize}
ℹtem Clue: +1 die on an investigation roll
ℹtem Small String: A minor favor or piece of information
ℹtem DV −1: On one roll within the informant's sphere of influence
\end{itemize}

\noindent\textbf{Informant Complications}
When an informant's Reliability reaches 0:
\begin{itemize}
ℹtem They flip allegiance or vanish
ℹtem Start Burn [4] timer to recover/replace
ℹtem GM may spend SB to create immediate problems
\end{itemize}

\noindent\textbf{Cross−Check Advantage}
Spend 2 informants on the same topic:
\begin{itemize}
ℹtem Take +1 Effect on related rolls
ℹtem Only the lowest Reliability decreases on a roll of 1
\end{itemize}

\noindent\textbf{Masking Source}
Spend 1 Bank to prevent Reliability −1 once when using informants.

\chapter{Court−Approved Heists (The Heist Timer )}

When the players accept a job from the Court, the GM sets up a \textbf{Heist Timer } [6] (simple) to [12] (complex) to track overall progress toward the objective. Additionally, the GM may establish \textbf{Threat Timers} to track escalating danger during the heist.

\chapter{Court−Approved Heists (The Heist Timer )}

\begin{heistengine}
\textbf{Heist Engine – GM Quick Reference}
\begin{tabular}{>{\raggedright\arraybackslash}p{0.32\textwidth} >{¢ering\arraybackslash}p{0.32\textwidth} >{\raggedleft\arraybackslash}p{0.32\textwidth}}
⊤rule
\textbf{If you want…} & \textbf{Do this…} & \textbf{The this…} \\
∣rule
\textbf{Tension to rise steadily} & Use 2−3 Threat Timers (Alarm, Patrol, Rival Interference). Tick them on \textit{Partials and Misses}. & Using only a single Heist Timer – players need to see \textit{multiple} pressures. \\
\textbf{Lockpicking to feel dramatic} & Make the safe a \textit{Progress Timer [6]}. Each Mechanics roll ticks \textit{1−2 segments on a Partial} (success on full tick). & Rolling once for "success/fail" – binary feels flat. \\
\textbf{Setup rolls to matter} & Grant \textit{Kit Tokens} (one−use +1d or Position boost) on \textit{strong hits}. On partials, start a complication timer (e.g., Guard Alert [3]). & Letting setup rolls have no mechanical effect. \\
\textbf{Boons to flow (not hoard)} & Remind players: \textit{"You have [X] Boons – spend them or lose them at scene end."} Model spending as GM by offering tempting re−rolls (e.g., "Spend a Boon to avoid that guard’s glance?"). & Letting players sit on multiple Boons for scenes – they’re meant to be spent. \\
\textbf{Flashbacks (preparedness)} & Let players spend \textit{1 Boon} to declare a minor flashback ("I loosened that window last night"). For major prep (gaining +2 dice), require a \textit{Setup roll} \textit{before} the heist. & Saying "no" to creative flashbacks – they’re the soul of heists. \\
\textbf{Position to affect stealth} & \textbf{State Position before every roll}:\\
& • \textit{Controlled} – "You have cover in the shadows."\\
& • \textit{Risky} – "The guard glances your way – act fast."\\
& • \textit{Desperate} – "The guard’s lantern sweeps \textit{right} at you." & Using Controlled as default – Position \textit{must} drive narration. \\
\textbf{Partials to feel like progress} & Offer a \textit{menu}: reduce next attempt’s DV by 1, improve Position by 1 step, tick a Progress Timer , gain a Clue, or \textit{succeed at a cost} (GM gains 1 SB). & Saying "nothing happens" – that’s a Miss (no effect, GM gains SB). \\
\textbf{Momentum after a Miss} & \textbf{Spend the SB immediately}: new guard appears, door locks/jams, or tick a Threat Timer \textit{twice}. & Banking SB – strike while the iron is hot (escalation feels punitive, not reactive). \\
\textbf{Teamwork to shine} & Remind players of \textit{Strings}:\\
& • \textit{"You and Benjamin share a String – invoke it now: you both gain 1 Boon or reduce DV by 1."}\\
& • Strings represent specific favors/debts (not vague "bonds"). & Letting players forget their Strings – they’re mechanical leverage. \\
\textbf{The heist to end memorably} & End with an \textit{escape roll} (even with \textit{Always Another Exit}). Failure means: dropped item, left clue, or gained \textit{Heat} (not just "you escape"). & Letting them walk out free – the last roll \textit{must} carry consequence. \\
⊥tomrule
\end{tabular}
\end{heistengine}

\section{The Velvet Heist Planner}
Before a job, each player describes one preparation (casing, bribing, stashing gear). Each strong success on a Setup roll grants +1 Position or −1 DV later. Each player starts with 2 \textbf{Kit Tokens} (lockpicks, dark lantern, false papers) that can be spent during the heist for:
\begin{itemize}
ℹtem +1 die on a roll
ℹtem Improve Position by one step
ℹtem Tick the Heist Timer forward by one segment
ℹtem Pause a Threat Timer (prevent it from advancing for one tick)
\end{itemize}
After use, the token is consumed.

\section{Heist Timer }
The Heist Timer advances when players:
\begin{itemize}
ℹtem Complete a phase of the heist (e.g., casing, infiltration, acquisition, escape)
ℹtem Acquire a key asset (a map, a key, a guard schedule)
ℹtem Eliminate a major obstacle (bribe a captain, disable an alarm)
\end{itemize}
The timer ticks backward when:
\begin{itemize}
ℹtem A guard raises an alarm
ℹtem A rival thief interferes
ℹtem The party is forced to leave evidence
\end{itemize}
When the timer fills, the heist is complete. The players gain the loot, plus 1 Favour with the Creditor who hired them. The GM may also award a \textit{Court Favour} (a one−use String that can be traded for help, information, or a safe house).
If the timer empties (reaches 0), the heist fails. The party must escape, may lose a member, and gains a \textit{Debt} owed to the Creditor (they will call it in later).

\section{Threat Timers}
The GM sets up Threat Timers to track specific dangers:
\begin{itemize}
ℹtem \textbf{Alarm [4]} − Security response time and alert level
ℹtem \textbf{Security [4]} − Methodical search patterns and guard rotations
ℹtem \textbf{Inside Trouble [4]} − Unexpected problems from within the target location
\end{itemize}
When a Threat Timer fills, the corresponding complication occurs (e.g., Alarm filling means guards are fully alerted and responding).


% −−−−−−−−−−−−−−−−−−−−−−−−−−−−−−−−−−−−−−−−−−−−−−−−−−−−−−−−−−−−−
% PART 5: PATRONS OF THE UNDERWORLD (Expanded)
% −−−−−−−−−−−−−−−−−−−−−−−−−−−−−−−−−−−−−−−−−−−−−−−−−−−−−−−−−−−−−
\part{Patrons of the Underworld}

\chapter{Ikasha −− She Who Sleeps in Shadow}

\section{Lore}
Ikasha is the hush between footfalls, the patience of dark water, the black−feathered watcher at every threshold. She does not command. She \emph{sleeps}. But in her dreams, she whispers possibilities.

Thieves who follow Ikasha learn to move unseen, to speak unremembered, to become shadows that slip between what is and what could be.

\section{Advantages for Rogues}
\begin{itemize}
ℹtem \textbf{Soft Witness} −− Your first Stealth roll each scene is Position +1.
ℹtem \textbf{Shadow Meld} −− Once per session, step into a shadow and emerge from another within Near range.
ℹtem \textbf{The Unobserved} −− When you are Hidden, Social SB cannot be used to frame you (unless you act first).
\end{itemize}

\section{Drawbacks}
\begin{itemize}
ℹtem \textbf{The Witness Debt} −− Every time you use Ikasha's power, you mark a \textit{Witness Timer } [4]. When it fills, a raven appears to watch your next crime. If you fail that roll, the raven reports to a rival.
ℹtem \textbf{Light Sensitivity} −− In bright conditions (direct sun, magical illumination), your Stealth rolls suffer −1 die.
\end{itemize}

\section{Rites for the Rogue Cantor}
\begin{itemize}
ℹtem \textit{The Unlatched Window} (Low) −− Gain +2 dice to open one lock or window silently.
ℹtem \textit{Raven Courier} (Standard) −− Send a whispered message up to a mile away via a spectral raven.
ℹtem \textit{Threshold of Sleep} (High) −− Walk through a locked door as if it were mist. Costs 2 Fatigue and marks Obligation.
\end{itemize}

\chapter{Livaea −− The Crimson Courtier}

\section{Lore}
Livaea is the whisper behind the curtain, the smile that sells secrets, the velvet hand that closes a cage. She rules through allure, performance, and the hunger for connection. Her domain encompasses not just courtiers, but also escorts, concubines, companions, and all who trade in intimacy as currency.

Thieves who follow Livaea do not steal coins−−−they steal \emph{trust}. A stolen kiss, a forged letter, a lover's secret sold to the highest bidder. Livaea's temples are the pleasure houses of Silkstrand, the massage parlours of Zakov, the discreet companionship agencies of the Brass Coast. Her priests are courtesans who know that a whispered secret is worth more than a night's fee, and that a carefully cultivated attachment can move armies.

\section{Advantages for Rogues}
\begin{itemize}
ℹtem \textbf{Courtly Graces} −− You gain +1 die to Sway and Performance in any refined setting.
ℹtem \textbf{The Velvet Glove} −− Once per session, you may declare a minor NPC is your secret admirer or blackmail victim. They will help you once.
ℹtem \textbf{Unreadable Mask} −− You may spend 1 Boon to ignore one Social SB spend against you (the scandal simply fails to stick).
ℹtem \textbf{Favoured Companion} −− You have a standing introduction to any pleasure house, escort agency, or discreet social club in the Amaranthine. The staff will not betray you unless the price is overwhelming.
\end{itemize}

\section{Drawbacks}
\begin{itemize}
ℹtem \textbf{The Jealous Court} −− When you use Livaea's power publicly, mark a \textit{Rival Timer } [4]. When it fills, a former paramour or jilted client becomes a recurring enemy.
ℹtem \textbf{Addicted to Adoration} −− If you go more than a session without a notable social triumph, you suffer −1 die to all social rolls until you perform a grand gesture (a public seduction, a scandalous revelation, a purchased favour).
ℹtem \textbf{The Velvet Price} −− Livaea expects her followers to keep the oldest profession thriving. Refusing to trade in intimacy (even as a broker, not a direct participant) may cause her to withdraw her favour.
\end{itemize}

\section{Rites for the Rogue Cantor}
\begin{itemize}
ℹtem \textit{Silk Whisper} (Low) −− Send a telepathic sentence to a willing target within sight.
ℹtem \textit{The Shared Cup} (Standard) −− Share a drink with a target; they cannot lie to you for the next hour.
ℹtem \textit{The Velvet Invitation} (Standard) −− Gain automatic, unquestioned entry to any private social function (masquerade, salon, pleasure house, nobles' retreat). The invitation appears in your pocket, signed by a name you can borrow for one night. Costs 1 Obligation.
ℹtem \textit{The Crimson Masquerade} (High) −− Assume the appearance of any person you have touched (skin−to−skin) for one scene. Your voice matches theirs perfectly. Costs 2 Fatigue.
ℹtem \textit{The Unrefusable Offer} (High) −− Sing a bargain while touching the target. They must accept or suffer −2 dice to all rolls until they do. The offer must be something they could conceivably agree to (no suicide, no self−mutilation). Costs 2 Obligation.
\end{itemize}

\chapter{Aveh −− The Rider Behind the Storm (For Fugitives)}

\section{Lore}
Aveh is the faceless rider at the horizon's edge. She offers freedom at the cost of belonging. Thieves who need to vanish−−−from a bad debt, a murder charge, a betrayed partner−−−pray to Aveh.

\section{Advantages}
\begin{itemize}
ℹtem \textbf{The Escape Clause} −− Once per session, you may declare an escape route that was not there before. The GM must accept it, but may impose a cost (you left something behind).
ℹtem \textbf{Forgotten Face} −− Witnesses have −2 dice to recall your description.
\end{itemize}

\section{Drawbacks}
\begin{itemize}
ℹtem \textbf{The Erased Name} −− Each time you use Aveh's power, you lose a minor personal connection (a contact forgets you, a safe house is no longer safe).
ℹtem \textbf{No Home} −− You cannot stay in one city for more than a season without attracting the Rider's attention.
\end{itemize}

\chapter{Malachai −− The Chained Angel (For the Desperate)}

See the expansion \textit{Malachai, the Chained Angel}. For thieves, she offers cursed luck−−−a die that always rolls high, but at the cost of a slow corruption. Use sparingly; Malachai is a patron for characters who are already falling.

\chapter{Choosing Your Patron}

\begin{tabular}{p{0.25\textwidth} p{0.35\textwidth} p{0.35\textwidth}}
⊤rule
Patron & Best for & Worst for \\
∣rule
Ikasha & Silent infiltration, ghost thieves & Social scenes, bright areas \\
Livaea & Courtesans, con artists, spies & Those who hate attention, the chaste \\
Aveh & Fugitives, escape artists & Anyone who wants roots \\
Malachai & Desperate gamblers, cursed builds & Long−term campaigns \\
⊥tomrule
\end{tabular}

% −−−−−−−−−−−−−−−−−−−−−−−−−−−−−−−−−−−−−−−−−−−−−−−−−−−−−−−−−−−−−
% APPENDICES
% −−−−−−−−−−−−−−−−−−−−−−−−−−−−−−−−−−−−−−−−−−−−−−−−−−−−−−−−−−−−−
\appendix
\part{Appendices}

\chapter{Sample Schemes}

\subsection{The False Heir}
Forge a letter of introduction from a distant relative. Infiltrate a noble house. Steal the heirloom. Disappear before the real relative arrives.  
\textbf{Timers:} Forgery [6], Infiltration [4], Escape [4].

\subsection{The Ghost Cargo}
Steal a shipment of silk. Replace it with crates of sand. Sell the silk in another port using forged manifests. The merchant spends months investigating "thieves" who never existed.  
\textbf{Timers:} Heist [6], Laundry [6], Investigation [8] (enemy).

\subsection{The Rival's Downfall}
Plant forged evidence of embezzlement in a rival thief's safehouse. Tip off the Watch. Claim their territory when they are arrested.  
\textbf{Timers:} Forgery [4], Infiltration [4], Frame Job [6], Territory Grab [4].

\chapter{NPC Profiles}

\subsection{The Watch Captain (Corrupt)}
\textbf{Name:} Captain Vellis.  
\textbf{Tier:} III.  
\textbf{Strings:} Patrol routes, jail keys, informant network.  
\textbf{Leverage:} He owes you for not exposing his affair.  
\textbf{Price:} 1 Favor = one night of ignored patrols.

\subsection{The Fence}
\textbf{Name:} Old Kes.  
\textbf{Tier:} II.  
\textbf{Strings:} Underground market, smuggler boats, counterfeit coin.  
\textbf{Price:} ▒ of the value, no questions. For an extra 10%, he will also launder the goods (start Laundry Timer at +1).

\subsection{The Rival Thief}
\textbf{Name:} "Quick" Lena.  
\textbf{Tier:} III.  
\textbf{Strings:} A network of urchin spies, a stolen set of master keys, a grudge against you.  
\textbf{Weakness:} She cannot resist a taunt. Use Deception to draw her into a trap.

\subsection{Court Informant}
\textbf{Name:} Scribe Marlo.  
\textbf{Tier:} I.  
\textbf{Strings:} Access to the Archivolt's night shift, a gambling debt.  
\textbf{Price:} 1 Boon = a copy of tomorrow's court docket (reveals who is being investigated).

\subsection{Livaean Courtesan (Information Broker)}
\textbf{Name:} Celeste.  
\textbf{Tier:} II (but with access to Tier IV secrets).  
\textbf{Strings:} A network of wealthy clients, a black book of secrets, a standing invitation to every salon.  
\textbf{Price:} A secret for a secret, or a favour for a night's companionship. She never sleeps with anyone she does not choose.

\chapter{Velvet Court Rumours (d12)}
Use these to seed adventures or colour scenes.

\begin{enumerate}
ℹtem The Grey Eminence is actually a Veiled Sister who defected. She knows the true name of a bound spirit that guards the Sunken Quarter.
ℹtem Old Kes is dying. He has hidden his fortune in three different locations. Whoever finds it first will inherit his smuggling network.
ℹtem Madam Serafine has a daughter. No one knows who the father is. The daughter is being raised by Sister Agatha.
ℹtem The Archivolt has a secret room containing every forged seal the Court has ever used. The key is held by the Archivist's ghost.
ℹtem The Dockmaster's warehouse on Pier Nine is a front for a Tulkani story−teller who trades in forgotten names.
ℹtem A new faction, the "Silk Purse," is trying to muscle in on the Velvet Court's territory. They are backed by a Kahfagian privateer.
ℹtem The Blue Ledger has a missing page. That page lists the name of a person who cannot be bribed.
ℹtem The Maze of Whispers connects to the Ways Between. Something comes through at midnight.
ℹtem The Ninth Key legend is real. One fragment is held by the Grey Eminence. Another is in the Archivolt's vault.
ℹtem Sister Agatha was once a Veiled Sister. She left the order after being ordered to bind a child's spirit.
ℹtem The Watch Captain's daughter is a member of the Velvet Court. She is the one who feeds him false information.
ℹtem The Velvet Court was founded by a Trow who wanted to study human greed. The Trow is still alive, waiting in the Under−arcades.
\end{enumerate}

\subsection{Rival Crew Generator}
When the GM needs a rival crew, roll on the tables below or choose options that fit the story.

\noindent\textbf{Name (d6)}
\begin{enumerate}
ℹtem Ashen Cord
ℹtem Brass Legion
ℹtem Thorns of Winter
ℹtem Silken Veil
ℹtem Iron Maw
ℹtem Shadow Consortium
\end{enumerate}

\noindent\textbf{Composition (Leader + 2−4 roles)}
Roll d6 for number of additional roles (2−5), then choose or roll for each:
\begin{itemize}
ℹtem Leader: The planner and face of the crew
ℹtem Enforcer: Handles violence and intimidation
ℹtem Fence: Moves stolen goods
ℹtem Informant: Gathers street intelligence
ℹtem Technician: Handles locks, traps, and devices
ℹtem Driver: Manages vehicles and escape routes
ℹtem Grifter: Specializes in cons and social manipulation
ℹtem Scout: Handles reconnaissance and surveillance
\end{itemize}

\noindent\textbf{Patron (d6)}
\begin{enumerate}
ℹtem A noble seeking deniable assets
ℹtem A merchant house protecting trade interests
ℹtem A temple faction advancing their agenda
ℹtem A foreign power gathering intelligence
ℹtem A criminal syndicate expanding influence
ℹtem A mysterious benefactor with unknown motives
\end{enumerate}

\noindent\textbf{Goal (d6)}
\begin{enumerate}
ℹtem Same artifact or information the PCs seek
ℹtem Eliminate the PCs to remove competition
ℹtem Claim a specific turf node or racket
ℹtem Expose the PCs to the authorities
ℹtem Disrupt a specific Velvet Court operation
ℹtem Accumulate wealth and influence in the underworld
\end{enumerate}

\noindent\textbf{Relationship (d6)}
\begin{enumerate}
ℹtem Open hostility − active conflict
ℹtem Professional rivalry − compete for same scores
ℹtem Uneasy truce − cooperate when beneficial
ℹtem Former allies − bad blood from past job
ℹtem Mentor−struck − one trained the other
ℹtem Family ties − complicated personal history
\end{enumerate}

\noindent\textbf{Resource (d6)}
\begin{enumerate}
ℹtem Specialized equipment (e.g., wardbreaker tools)
ℹtem Network of informants in a specific district
ℹtem Access to legitimate business fronts
ℹtem Political connections in the city guard
ℹtem Magical or supernatural assistance
ℹtem Large cash reserves for bribes and gear
\end{enumerate}

\noindent\textbf{Agenda Timer [6]}
Track the rival crew's progress toward their goal. Tick the timer when:
\begin{itemize}
ℹtem The party ignores them or their activities
ℹtem The GM spends 2+ SB on their off−screen activity
ℹtem They successfully complete a step toward their goal
\end{itemize}
When the Agenda Timer fills, the rival crew achieves their goal (which may complicate the PCs' plans).

% Appendix D: The Silk House — A Livaean Brothel Asset
% Insert into your Fate's Edge document's appendix section.

\chapter{The Silk House — A Livaean Brothel Asset}
łabel{app:silk−house}

\begin{quote}
ℹtshape
A well−run house of pleasure is a fortress of secrets. The coin paid at the door buys only the first hour. The second hour costs a confession. The third buys a silence that can never be broken.

\vspace{0.2cm}
−−− Madam Serafine, \textit{The Velvet Touch}
\end{quote}

\section{The Velvet Altar}

In the back rooms of Silkstrand’s finest pleasure houses, behind curtains dyed with madder and saffron, small shrines to Livaea stand in alcoves. A candle, a mirror, a single silk ribbon tied in a lover’s knot. Livaea does not demand celibacy or shame. She blesses the transaction: the honest exchange of comfort for coin, of secrecy for safety, of power for pleasure.

A brothel dedicated to Livaea is not merely a business. It is a \textbf{threshold} — a place where masks are laid aside, where the powerful become vulnerable, where a whisper in the dark can ruin a dynasty. The courtesans are priestesses. The Madame is her high priestess. And every client who crosses the threshold leaves behind a piece of themselves — a secret, a preference, a name that should not have been spoken.

The Velvet Court has long maintained an understanding with Livaea’s faithful. Information flows both ways. A thief who insults a Livaean courtesan finds every door closed. A patron who pays generously may find a lock unpicked by morning.

\subsection{Running a Velvet Racket}
When the party controls a territory (e.g., the Dye District docks, the Three−Queens Bridge approach), the GM notes this as a Turf Node. Each node has a Racket (pick from: \textit{Fencing Operation, Smuggling Cut, Protection Circle, Information Stall}). Holding turf requires a Hold roll (Wits+Tactics vs DV 2−4) each session. On success, gain Bank +1 or local goodwill. On failure, rivals contest the node—start a War [8] timer.

\noindent\textbf{Example Turf Node:}
\textit{Dye District Docks} − Racket: Fencing Operation

\noindent\textbf{Hold Roll Outcomes:}
\begin{itemize}
ℹtem \textbf{Hit:} Bank +1 or local goodwill (reduce DV by 1 on one social roll made in this territory)
ℹtem \textbf{Miss:} War [8] timer starts with 1 segment filled (rivals begin contesting control)
\end{itemize}


\section{The Silk House as an Asset (Major Asset, 12 XP)}

A Silk House is a \textbf{Major Asset} (12 XP) that provides off‑screen leverage through blackmail, rumour, and social access. It cannot be purchased quickly; establishment requires one month of in‑game time, during which the party must recruit staff, secure premises, and build a clientele (or take over an existing establishment).

\begin{description}
  ℹtem[Cost] 12 XP
  ℹtem[Establishment] 1 month (or as story permits)
  ℹtem[Upkeep] Pay 12 XP per downtime \emph{or} devote a full scene to managing it (settling disputes, bribing watch, maintaining cover).
  ℹtem[Capacity (Optional)] Counts as one Asset against your Spirit limit.
\end{description}

\subsection{Free Use (Once per Session, Off‑Screen)}

Choose \textbf{one} of the following benefits at the start of a scene. The effect is narrative, not rolled.

\begin{description}
  ℹtem[Blackmail File] You have compromising information on a named minor NPC (guard captain, customs clerk, merchant). Once, you may use this file to \textbf{improve Position by one step} on a social action against that NPC or to \textbf{demand a minor favour} (e.g., look the other way, delay a patrol). The file is consumed after use; you must spend a downtime scene gathering new material (part of upkeep).

  ℹtem[Rumour Mill] Before a scene set in the city, you may state a plausible rumour about a person, place, or upcoming event. The GM must incorporate it as a truth (though possibly incomplete or misleading). Examples: \textit{“The Watch Captain visits the Silk House every Restday.”} or \textit{“The Dye Syndicate is about to raise prices.”}

  ℹtem[Soft Influence] You may skip one social gatekeeping roll (e.g., gaining an audience, obtaining an invitation) by stating that a client of the Silk House made the introduction. The GM may ask for a name or a short flashback, but no roll is required.

  ℹtem[Safe Harbour] A fugitive or witness can hide in the Silk House for one scene without risk of discovery (unless the searchers explicitly search brothels). Use this to extract a contact or to plan a heist.
\end{description}

\subsection{Scene Surge (Spend 1 Boon, On‑Screen)}

Once per session, you may spend \textbf{1 Boon} to trigger one of the following on‑screen interventions:

\begin{description}
  ℹtem[The Velvet Curtain] A courtesan distracts a guard or a mark at a critical moment. The next Stealth or Subterfuge roll in the scene gains \textbf{+1 die}.

  ℹtem[A Whisper in the Right Ear] An influential client intervenes on your behalf. Reduce \textbf{DV by 1} on a single Petition or Broker roll.

  ℹtem[The Back Stair] You escape through a hidden passage known only to the house. Convert a pursuit consequence into a temporary complication (e.g., you drop a clue instead of being caught).

  ℹtem[A Taste of Honey] A courtesan serves a target a drink laced with a mild truth serum (or simply loosens their tongue). The next time you roll Sway against that target, treat \textbf{Position as Dominant}.
\end{description}

\subsection{Extra Use}

You may spend \textbf{2 XP} to trigger the Free Use benefit an additional time this session.

\section{The Madame as a Terrestrial Patron}

The Silk House is not just an asset; it is a \textbf{living organisation} with its own will, its own debts, and its own demands. The Madame (named NPC, Tier II or III) acts as a terrestrial patron. Characters who wish to draw deeper on the House’s resources may swear an oath to Livaea (or to the Madame) and incur \textbf{Obligation}.

\subsection{Obligation Track}

\textbf{Capacity}: Spirit + 2.  
When you call on the House’s resources beyond the normal Free Use, mark \textbf{Obligation} as follows:

\begin{center}
\begin{tabular}{ll}
⊤rule
\textbf{Request} & \textbf{Obligation Cost} \\
∣rule
Use Free Use a second time in a session & +1 \\
Ask the Madame to hide a fugitive for a week & +2 \\
Request a forged letter of introduction to a high noble & +2 \\
Borrow a Livaean Rite (e.g., \textit{Golden Tongue} or \textit{The Unrefusable Offer}) & +3 \\
Demand the House use its blackmail to ruin a rival (major action) & +4 \\
⊥tomrule
\end{tabular}
\end{center}

When \textbf{Obligation exceeds Capacity}, the Madame sends a polite request: a favour to be named later. If you refuse, the House’s Free Use is \textbf{unavailable} until you clear the debt. If you ignore her twice, she sells one of your secrets to a rival (GM chooses a String you lose).

\subsection{Madame’s Favours (Clearing Obligation)}

You may reduce Obligation by:

\begin{itemize}
  ℹtem \textbf{Performing a Service} (e.g., eliminate a rival brothel, intimidate a blackmail victim, deliver a secret to a client) – reduces Obligation by \textbf{2–4} depending on difficulty.
  ℹtem \textbf{Paying Coin} – 1 Boon reduces Obligation by \textbf{1}; 1 Diamond (e.g., a Forged Document or valuable item) reduces Obligation by \textbf{2}.
  ℹtem \textbf{Sponsoring a New Courtesan} – Recruit a talented individual (Cap 1–2 follower) to work at the House. Reduces Obligation by \textbf{3}.
\end{itemize}

\subsection{Sample Madame: Madam Serafine’s Rival}

\begin{description}
  ℹtem[Name] \textbf{Madam Illara} (Tier III)
  ℹtem[Motivation] Expand her influence into the Brass Coast; compete with Serafine’s Velvet Court.
  ℹtem[Quirk] She lost an eye to a jealous client; wears a ruby in the socket that glows faintly when someone lies.
  ℹtem[Strings] Clients among the Admiralty, the Dye Syndicate, and the Archivolt.
  ℹtem[Leverage] She knows which notary can be bought for a night’s company.
\end{description}

\section{Upkeep and Complications}

\subsection{Maintaining the House}

Each downtime, you must perform \textbf{one} of the following:

\begin{itemize}
  ℹtem Pay \textbf{12 XP} (the Asset’s base cost) – represents bribes, repairs, and salaries.
  ℹtem Devote a \textbf{full scene} to managing the House: settling disputes among staff, bribing the Watch, smoothing over a scandal, or acquiring new blackmail material.
  ℹtem Perform a \textbf{Service for Livaea} (GM’s choice) – e.g., sponsor a festival, sabotage a purity crusade, deliver a secret to a Livaean shrine.
\end{itemize}

If you neglect upkeep for \textbf{one} downtime, the House becomes \textbf{Neglected}:
\begin{itemize}
  ℹtem Free Use only works half the time (GM’s discretion; roll 1d6, 1–3 = fail).
  ℹtem Obligation cannot be cleared by Boons (only by services).
\end{itemize}

If neglected for \textbf{two} consecutive downtimes, the House becomes \textbf{Compromised}:
\begin{itemize}
  ℹtem The Asset provides \textbf{no} Free Use until restored.
  ℹtem A rival brothel gains +1 Favour with the Watch.
  ℹtem The GM may advance a \textbf{Rival Timer } (e.g., \textit{Watch Raid} or \textit{Syndicate Takeover}).
\end{itemize}

Restoring a Compromised House requires a \textbf{full arc} (2–3 sessions) of in‑game work, not just XP payment.

\subsection{Complications (d6)}

At the start of any scene where you use the Silk House’s Free Use or Scene Surge, the GM may roll (or choose) a complication:

\begin{enumerate}
  ℹtem \textbf{A Jealous Rival} – Another pleasure house sends thugs to intimidate your staff. Mark +1 Obligation to the Madame or lose 1 Boon.
  ℹtem \textbf{Watch Inspection} – A corrupt captain demands a bribe. Pay 1 Boon or the House becomes \textit{Neglected} for this session.
  ℹtem \textbf{A Client’s Secret Leaks} – Someone you blackmailed threatens to expose you. Gain Exposure +1 (city faction).
  ℹtem \textbf{Livaea’s Price} – The goddess demands a favour: sponsor a festival, deliver a cursed mirror, etc. Refuse and the House’s Free Use is unavailable for one session.
  ℹtem \textbf{Purity Crusade} – A temple zealot preaches against the House. Social rolls in that district are –1 Position until you bribe or discredit them.
  ℹtem \textbf{The Hollow’s Interest} – A Oath‑Keeper is seen loitering outside. The House is safe, but the party feels watched. The GM gains 1 SB to spend on a future complication.
\end{enumerate}

\section{Sample Blackmail Files (d6)}

When you use the \textit{Blackmail File} Free Use, roll or choose what you have:

\begin{enumerate}
  ℹtem \textbf{The Watch Captain’s Affair} – He visits every Restday. Leverage: one ignored patrol.
  ℹtem \textbf{The Merchant’s False Ledger} – You have a copy of his tax fraud. Leverage: 20% discount on his goods.
  ℹtem \textbf{The Notary’s Forged Seal} – He sealed a treaty that never existed. Leverage: one document stamped without question.
  ℹtem \textbf{The Admiral’s Illegitimate Child} – He pays for the child’s silence. Leverage: safe passage through harbour patrols.
  ℹtem \textbf{The Priest’s Coded Letters} – He corresponds with a rival temple. Leverage: he will preach a sermon of your choosing.
  ℹtem \textbf{The Noble’s Secret Marriage} – He married below his station. Leverage: one invitation to a closed salon.
\end{enumerate}

\section{Adventure Hooks}

\begin{itemize}
  ℹtem \textbf{The Purge} – A new temple faction demands the closure of all Livaean houses. The party must defend the Silk House – legally, violently, or by exposing the zealot’s own hypocrisy.

  ℹtem \textbf{The Black Book} – A former client steals the House’s blackmail ledger. He threatens to publish it unless the party retrieves a cursed artifact for him.

  ℹtem \textbf{The Rival Courtesan} – A newly arrived courtesan is actually a spy for the Veiled Sisters. She is using the House to identify Livaean followers for elimination.

  ℹtem \textbf{The Ninth Client} – A stranger pays for a night but never leaves his room. He is a Oath‑Keeper, waiting for the House to offer him a secret as tribute. The party must learn what he wants before he collects.
\end{itemize}

\section{Integration with Velvet Court and Livaea’s Cult}

If the party already belongs to the \textbf{Velvet Court} (see \textit{The Velvet Touch}), the Silk House can be a \textbf{Creditor’s asset}. Madam Serafine may own a share. The party’s Velvet Dial may affect House benefits (e.g., at +2 Dial, the Free Use ignores one complication per session). At –2 Dial, the House is raided by the Watch.

For followers of \textbf{Livaea}, the Silk House counts as a \textbf{Minor Shrine}. You may learn her Rites there (e.g., \textit{Golden Tongue}, \textit{The Unrefusable Offer}) without needing a separate temple. Using a Livaean Rite within the House reduces its Obligation cost by 1 (minimum 1).

\begin{quote}
ℹtshape
The silk sheets hide more than skin. They hide the ledger of the soul. Treat the House well, and it will pay you back a thousand times. Treat it poorly, and you will find your own name on a very different kind of list.

\vspace{0.2cm}
−−− Madam Serafine, \textit{The Velvet Touch}
\end{quote}

\end{document}% This is a fragment; do not include a full preamble when inserting into a larger document.
\chapter*{Colophon}
\addcontentsline{toc}{chapter}{Colophon}
This book was written in a back room of the Velvet Court, on paper that had once been a customs manifest. The ink is lamp−black cut with the ash of a burned ledger. The binding is silk thread, dyed red with madder root and a single drop of my own blood−−−for luck.

The lantern gutters. The night is young. There is a mark waiting at the Counting House, and he does not yet know that his fortune is already spent.

\begin{flushright}
\textit{Madam Serafine}\\
\textit{Silkstrand, under the Three−Queens Bridge}\\
\textit{The ink is fresh. The coin is warm.}
\end{flushright}

\end{document}